\section{Introduction}
Consider the problem $\holant{=_2}{\mathcal{F}}$, where $\mathcal{F}$ is closed under conjugation. i.e., $f\in \mathcal{F}\iff \overline{f}\in \mathcal{F}.$

\begin{theorem}\label{thm: tractable classes}
    Let $\mathcal{F}$ be a set of complex valued signatures.
    Then $\Holant(\mathcal{F})$ is tractable if
    \begin{equation}\label{cond: tractable classes}
    \mathcal{F}\subseteq \mathscr{T}, ~~~
     \mathcal{F} \text{ is }\mathscr{P}\text{-transformable,}~~~
     \mathcal{F} \text{ is } \mathscr{A}\text{-transformable,~~~or~} 
     \mathcal{F} \text{ is } \mathscr{L}\text{-transformable.}\tag{$\mathrm{ \textcolor{red}{T}}$}
     \end{equation}
\end{theorem}

See \cref{subsec: tractable signatures} for their definitions.

\begin{theorem}[Main theorem]\label{thm: main theorem}
    Let $\mathcal{F}$ be a complex signature set that is closed under conjugation. 
    Then $\Holant(\mathcal{F})$ is \#P-hard unless $\mathcal{F}$ satisfies {\rm (\ref{cond: tractable classes})}, in which case it is tractable.
\end{theorem}

\paragraph{Proof Overview (Draft)}
    The tractable part in the main theorem follows from \cref{thm: tractable classes}.
    We only need to prove that $\Holant(\mathcal{F})$ is \#P-hard if $\mathcal{F}$ doesn't satisfy {\rm (\ref{cond: tractable classes})}.
    We assume ${\mathcal{F}}\not\subseteq \mathscr{T}$.
    Also, since  ${\mathcal{U}}^{\otimes}\subseteq \mathscr{T}$, ${\mathcal{F}}\not\subseteq {\mathcal{U}}^{\otimes}$.
    Thus, there is a nonzero signature ${f}\in {\mathcal{F}}$ of arity $2n$ such that ${f}\notin {\mathcal{U}}^{\otimes}$.
    We want to achieve a proof of  \#P-hardness by induction on $2n$. 
    The base cases $2n=2$ is handled by {\sc 1st-Orth} (\cref{lm: not 1st-orth is hard}), $2n=4$ is handled by {\sc 2nd-Orth} and non-existence of the AME(4,2) state (\cref{lm: base case 2n=4}).

    For higher arity case, our hope is to ``realize" $f'$ from $f$, with lower arity but keeping the property $f'\notin \mathcal{U}^\otimes.$
    Common arity reduction realization methods are merging, mating and factorization.
    However, this hope is an illusion, due to the extrodinary signatures $f_6$ and $f_8$ discovered in \cite{realholant}.
    Themselves are not in $\mathcal{U}^\otimes$, but all signatures realizable by merging, mating and factorization lie in $\mathcal{U}^{\otimes}.$
    That's why we must do individually tailored proofs for $2n=6$ and $2n=8$.

    For $2n=6$, our idea is to first realize $f_6$ (\cref{sec: Third Order Orth,sec:recover-f6}), then prove $\Holant(f_6,\mathcal{F})$ is $\#P$-hard under the assumption that $\mathcal{F}\not \subseteq \mathscr{A}$ (\cref{sec: f6 is hard}).
    To realize $f_6$, we first realize an AME(6,2) state (or equivalently, a signature satisfying {\sc 3rd-Orth}) in \cref{sec: Third Order Orth}.
    This utilizes the framework of Xia \cite{MingjiProgram}, where proved the ``projective binary group" generated by $f$ must form an finite subgroup of $\mathbf{PU}(2)=\mathbf{SU}(2)/\{\pm I_2\}$.
    We give an exposition of this proof in \autoref{appendix: mingji}.
    Since $\mathbf{SU}(2)/\{\pm I_2\}\cong \mathbf{SO}(3)$, and there is a finite list of finite subgroups of $\mathbf{SO}(3)$, we discuss all possible cases of the projective binary group in \cref{sec: Third Order Orth}.
    For the polyhedral group case (\cref{subsec: polyhedral groups}), large dihedral group case (\cref{subsec: dihedral group}) and standard case of Klein group (\cref{subsubsec: Standard Case Klein}), we prove that a signature satisfying {\sc 3-rd Orth} can be realized.
    We also prove the non-standard case of Klein group (\cref{subsubsec: Non-standard Klein}) and the cyclic group case (\cref{subsec: Trivial group,subsec: Order Two group,subsec: Cyclic group of Order at least three}) can not happen.
    Then we appeal to a result in quantum information theory, i.e. the uniqueness of the AME(6,2) state up to a basis change and permutation of wires (see \autoref{sec: AME62 has a common local unitary normal}). By this uniqueness, the realized signature is, up to a nonzero
    scalar, locally unitarily equivalent to $f_6$. However, the
    corresponding local unitary matrices may not be realized. In \cref{sec:recover-f6}, we show that, either $\#P$-hard, or a real orthogonal holographic transformation places these unitary matrices
    projectively in the tetrahedral group. We then prove that $f_6$ and four Bell signatures $\mathcal{B}$ can be realized by considering the two possibilities for the available projective binary group:
    it is either the tetrahedral group $\Gamma$ or a subgroup of $K_4$.
    

    In \cref{sec: f6 is hard}, we prove hardness when $f_6$ can be realized.
    The proof consists of two parts.
    If every signature in $\mathcal{F}$ has parity, we can do realification and incur the result in real Holant (\cref{subsec: f6 parity}).
    Otherwise, we can do arity reduction while preserving the non-parity property, until we realize a binary signature without parity.
    Then the proof is finished by extending the ``non-$\mathcal{B}$ hard" property of $f_6$ in \cite{realholant} to the ``non-$\Gamma$ hard" property, and a minimal counterexample argument (\cref{subsec: f6 non-parity}).

    For $2n=8$, we first realize $f_8$ (\cref{subsec: realize f8 Reed Muller}), and then prove $\Holant(f_8,\mathcal{F})$ is \#P-hard under the assumption that $\mathcal{F}\not \subseteq \mathscr{A}$ (\cref{subsec: f8 available hardness}).
    The proof follows the real Holant framework \cite{realholant}, and the realification method similar to \cref{subsec: f6 parity} again shows its power.

    For $2n\ge 10$, there are no more obstacles and we can do arity reduction while preserving the property $f\notin \mathcal{U}^{\otimes}$, thus completing the induction proof (\cref{sec: 2n >= 10}).